\textcolor{black}{Sebagaimana telah disinggung di \ref{chap:Dasar_Teori_Pendukung}, mahasiswa dituntut untuk mengevaluasi solusi-solusi yang mungkin berdasarkan proses dan standar keteknikan. Untuk mencapai poin tersebut, di Bab ini, mahasiswa harus memilih solusi terbaik untuk permasalahan spesifik yang telah dipilih.} Setelah dideskripsikan kekurangan dan kelebihan metode, tim akan memilih metode yang paling tepat disertai pengembangan / modifikasi / penyesuaian metode tersebut dengan mengacu pada pemodelan permasalahan yang telah ditulis pada \ref{chap:Pemodelan_Permasalahan} dan aspek - aspek relevan seperti yang telah ditulis pada Subbab~\ref{subsec:Penjabaran_Masalah_Umum_Menjadi_Masalah_Teknis} dan Subbab~\ref{subsec:Kendala-Kendala_yang_Perlu_Dipertimbangkan}. Jika dari semua metode yang ada tidak menghasilkan solusi yang konkret, tim bisa membuat metode sendiri. Dalam hal ini, proyek \textit{capstone} ini akan menjalankan mode penelitian sehingga tidak diwajibkan untuk menghasilkan luaran tertentu. Walaupun begitu adanya luaran seperti tertulis pada lembar luaran sangat disarankan.
\vspace{6pt}

\noindent
Contoh (sederhana):

"Model matematis larutan oksigen terlarut yang telah dirumuskan pada Persamaan~\ref{eqn:model_matematis_oksigen} cukup linear dan relatif sederhana. Akan tetapi, kondisi perairan payau di Indonesia cukup membawa tantangan seperti:
\begin{itemize}
    \item fluktuasi suhu dan salinitas (karena pengaruh pasang surut),
    \item gangguan konsumsi oksigen dari organisme seperti udang, ikan, dan mikroorganisme (yang bisa sangat bervariasi),
    \item masuknya bahan organik secara tiba-tiba (dari limbah atau hujan deras membawa material organik), dan
    \item gangguan lingkungan alami yang menyebabkan ketidakpastian.
\end{itemize}
Oleh karena itu, meskipun model awalnya sederhana, realitanya sistem mengalami dinamika tak pasti dan gangguan yang besar. Dari ketiga metode kendali potensial yang telah dibahas sebelumnya, FLC adalah metode yang paling sesuai untuk pengendalian kadar oksigen terlarut di perairan payau. Berikut adalah perbandingan ketiga metode tersebut yang disajikan pada Tabel~\ref{tab:Ch05_perbandingan_dan_pemilihan_metode_kendali}."
\begin{longtable}{|p{3cm}|p{3.5cm}|p{3.5cm}|p{3.5cm}|}
    \caption{Perbandingan Metode Pengendalian Kadar Oksigen Terlarut} 
    \label{tab:Ch05_perbandingan_dan_pemilihan_metode_kendali}
    \vspace{-.75em}\\
    \hline
     & \textbf{Pengendali PID} & \textbf{FLC} & \textbf{MPC} \\ \hline
    \centering \textbf{Kelebihan} & Simpel, cepat, murah diterapkan. & Tahan terhadap ketidakpastian, tidak perlu model presisi. & Optimal, bisa memprediksi masa depan dan mempertimbangkan banyak variabel. \\ \hline
    \centering \textbf{Kekurangan} & Sulit beradaptasi dengan perubahan sistem dan gangguan besar. & Perlu desain aturan \textit{fuzzy} yang hati-hati dan pengalaman operator. & Membutuhkan model akurat dan komputasi tinggi. \\ \hline
    \centering \textbf{Relevansi untuk Perairan Payau Indonesia} & Kurang optimal untuk kondisi perairan payau yang sangat dinamis. & Cocok, karena bisa menangani ketidakpastian dan perubahan lingkungan. & Terlalu kompleks dan berat untuk banyak tambak/perairan payau sederhana. \\ \hline
\end{longtable}
\vspace{12pt}

Pada contoh di atas, untuk tujuan penyederhanaan tulisan, aspek-aspek relevan yang lain masih belum dibahas. Namun, dalam praktiknya, Anda perlu mempertimbangkan aspek - aspek relevan seperti yang telah ditulis pada Subbab~\ref{subsec:Penjabaran_Masalah_Umum_Menjadi_Masalah_Teknis} dan Subbab~\ref{subsec:Kendala-Kendala_yang_Perlu_Dipertimbangkan} dalam memilih solusi yang tepat.

